\documentclass[10pt,a4paper]{article}
\usepackage[utf8]{inputenc}
\usepackage{amsmath}
\usepackage{amsfonts}
\usepackage{amssymb}
\usepackage{enumerate}
\title{GSF-3100 \\
02 Mathématiques financières des obligations (fraction de période) \\
Série d'exercices 2}

\begin{document}
\maketitle

\begin{itemize}
\item Dans ce document vous avez des questions à choix multiples allant de a) à d).  
\item Il y a seulement un choix possible par question.
\item Ce questionnaire contient des exercices en lien avec le principe de fraction de période.
\item Contient 12 questions. 
\end{itemize}
 

\newpage

\section*{Section 1}
Vous considérez l’achat d’une obligation versant des coupons annuels le 31 décembre et dont le dernier coupon a été versé il y a 260 jours.  En supposant que vous êtes dans une année non bissextile, calculer la fraction de période k nécessaire pour déterminer l’intérêt couru en utilisant la méthode:

\subsection*{1.  Exact / Exact} 
\begin{enumerate}[a)]
\item 0,74133
\item 0,70113
\item 0,71233
\item 0,72523
\end{enumerate}

\subsection*{2.  Exact/365} 
\begin{enumerate}[a)]
\item 0,71233
\item 0,71933
\item 0,70133
\item 0,72003
\end{enumerate}

\subsection*{3. Exact/360}
\begin{enumerate}[a)]
\item 0,71233
\item 0,72222
\item 0,71111
\item 0,72933
\end{enumerate}

\subsection*{4. 30/360}
\begin{enumerate}[a)]
\item 0,72933
\item 0,72233
\item 0,70389
\item 0,71389
\end{enumerate}

\newpage

\section*{Section 2}
Vous considérez l’achat d’une obligation versant des coupons les 30 juin et 31 décembre de chaque année. Calculer la fraction de période $k$ nécessaire pour déterminer l’intérêt couru en supposant que vous êtes le 4 février d’une année non bissextile et en utilisant la méthode:


\subsection*{5.  Exact / Exact} 
\begin{enumerate}[a)]
\item 0,18889
\item 0,19337
\item 0,19444
\item 0,19178
\end{enumerate}

\subsection*{6.  Exact/365} 
\begin{enumerate}[a)]
\item 0,18889
\item 0,19337
\item 0,19444
\item 0,19178
\end{enumerate}

\subsection*{7. Exact/360}
\begin{enumerate}[a)]
\item 0,18889
\item 0,19337
\item 0,19444
\item 0,19178
\end{enumerate}

\subsection*{8. 30/360}
\begin{enumerate}[a)]
\item 0,18889
\item 0,19337
\item 0,19444
\item 0,19178
\end{enumerate}

\newpage 

\section*{Section 3}
Vous considérez l’achat d’une obligation versant des coupons au dernier jour de chaque mois et dont le prochain coupon sera versé dans 7 jours. Calculer la fraction de période $k$ nécessaire pour déterminer l’intérêt couru en supposant que vous êtes en mars et en utilisant la méthode: 

\subsection*{9.  Exact / Exact} 
\begin{enumerate}[a)]
\item 0,80000 
\item 0,7800
\item 0,77419 
\item 0,78904 
\end{enumerate}

\subsection*{10.  Exact/365} 
\begin{enumerate}[a)]
\item 0,77419  
\item 0,78000
\item 0,80000 
\item 0,78904
\end{enumerate}

\subsection*{11. Exact/360}
\begin{enumerate}[a)]
\item 0,77419  
\item 0,78000 
\item 0,80000 
\item 0,78904
\end{enumerate}

\subsection*{12.  30/360}
\begin{enumerate}[a)]
\item 0,77419  
\item 0,78000 
\item 0,80000 
\item 0,78904
\end{enumerate}


\end{document}